\section{Complete Q support and effective matching dimension}
\label{sec:feasible-span}

The preceding generic Q phase makes a complete quadrilateral-corner list
available when positive Euler characteristic fails to produce an essential
disc. Report 81 already isolates forced-zero removal as an admission rule
for lower-dimensional geometry (Section~\ref{sec:planar-overlay-refinements}).
This continuation obtains the maximal feasible support from that completed
list, recompiles the source matching kernel, and uses the effective dimension
for standard enumeration. It also lets the independent checker verify this
support information before rebuilding a smaller standard cone.

\subsection{A completed ray list determines every forced-zero type}

For the original allowed support write
\[
 C=\{q\in\mathbb R^k:q\ge0,\ A_\Sigma q=0\}.
\]
Let $r_1,\ldots,r_m$ be its complete primitive Q extreme directions and set
\[
 J=\{j:\text{some }(r_i)_j>0\}.
\]
The set $J$ is computed by scanning all coordinates in the completed list.
Neither a successful early Q witness nor a capped or interrupted list
provides the required completeness premise.

\begin{proposition}[Maximal support from complete Q rays]
Every feasible quadrilateral vector vanishes off $J$. Restricting the
matching system to the coordinates in $J$ preserves the whole standard
normal cone after reinstating those zero coordinates, including all
non-link standard rays and their actual source coordinates. The restricted
matching kernel equals the linear span of its feasible Q cone.
\end{proposition}
\begin{proof}
The nonnegative matching cone is pointed. Its section $\sum_jq_j=1$ is
compact, so every nonzero feasible vector is a conic combination of its
extreme directions. Off-$J$ coordinates therefore vanish in every feasible
vector. Conversely $q_* = \sum_i r_i$ is feasible and strictly positive
on every coordinate in $J$. For any vector $u$ in the restricted matching
kernel, sufficiently small positive $\epsilon$ makes both
$q_*+\epsilon u$ and $q_*-\epsilon u$ nonnegative on $J$; their difference
shows that $u$ lies in the feasible cone's linear span.
Every feasible standard vector projects to a feasible Q vector, so its
removed coordinates were already zero. Conversely reinstating these zeros
in a feasible restricted standard vector preserves all original matching
and nonnegativity equations. The two standard cones are thus linearly
isomorphic. This isomorphism preserves extreme rays, including the vertex
links; canonical normalization and the full source coordinates of each
non-link ray also agree. No topology or component query is inferred from
positive Euler characteristic alone.
\end{proof}

This implements the maximal-support conclusion of the report's LP theorem
through a different available witness. It does not install a polynomial LP
solver. Obtaining the complete Q list still costs the preceding generic
Q search in the \emph{original} matching dimension. Recompilation costs are
charged rather than hidden in a dimension parameter.

\subsection{Adaptive recompilation after a completed positive Q phase}

At release \code{936e83365}, automatic supplied-sector discovery above
nullity three retains the Q-first policy described here. The sufficient
matching preprocessor of Section~\ref{sec:matching-support} later precedes
this complete-Q route and can avoid its original-dimensional enumeration. A verified disc returns immediately. A nonpositive complete
Q phase returns its ordinary Q exclusion proof. Only
\code{POSITIVE\_EULER\_ONLY} reaches the support step: all Q corners have
been tested, and none supplied an essential disc.

If $J$ is the full allowed support, the already prepared original kernel
is reused. Otherwise the implementation rebuilds the supplied source kernel
with exactly those allowed types in $J$. The ordinary automatic dispatcher
then selects the envelope, planar, or generic standard method using this
new kernel's actual nullity. No original coordinate is merely overwritten
with a guessed zero; matching geometry is reconstructed from the source.
The Q and standard phases retain their shared candidate allowance.
Exhaustion in either phase is inconclusive, and interruptions in support
scanning, recompilation or proof translation propagate.

Statistics describe the kernel that performed the standard search. A
\code{support\_reduction} record retains the requested kernel statistics,
original-coordinate indices and number of removed types. A positive witness
still contains the actual full triangulation coordinates and the original
allowed type list; its existing disc proof suffices independently of the
support optimization.

\subsection{A source proof that still states the original sector}

A negative standard result retains the existing
\code{normal-sector-exhaustion-v1} schema, original source digest, original
allowed types and standard phase. Its Q records are expanded to the original
coordinate width by inserting the removed zeros. The optional
\code{q\_support\_certificate} is the complete original Q-phase exhaustion
proof with status \code{POSITIVE\_EULER\_ONLY}. The parent never substitutes
a smaller requested sector for the original statement.

The independent checker first validates that the nested proof has the same
original source and allowed types, the Q phase and the required status.
Nested support annotations inside it are rejected, so there is only one
level of support proof. The checker then reconstructs and verifies the
complete original Q list by its own dense elimination and reference
enumeration. It derives $J$ from that authenticated list; it accepts no
producer-supplied retained-index list. Only then does it independently
rebuild and enumerate the restricted standard model, restore the original
coordinate width, and compare every parent ray, Euler value and negative
disc certificate against that model.

A malformed or incomplete nested Q proof cannot license removing any
coordinate. An old certificate without the annotation still follows the
original full-source checker. Removing a valid annotation also produces a
valid original full-sector certificate: its complete parent ray list and
source statement are unchanged. The optimization therefore accelerates
reconstruction rather than changing the negative statement's meaning.

\subsection{Raw and effective dimensions have different costs}

Let $d$ be the original matching nullity, $s$ the rebuilt nullity and $k_J$
the retained support size. Once the complete Q list is available, finding
$J$ costs $O(km)$ scalar inspections, with callbacks and bit access charged.
The rebuilt kernel has nullity exactly the feasible Q span dimension $s$.
Standard geometry consequently has the preceding
$\operatorname{poly}(T+B)(1+k_J)^{O(s)}$ enumeration bound. For $s\le3$ the
native low-dimensional routines apply; in particular the complete planar
source ray count is quadratic in retained support size.

The overall \emph{ray geometry} cost still includes the original Q phase,
whose upper bound is
$\operatorname{poly}(T+k+B)(1+k)^{d-1}$, plus support scanning, source
recompilation and the reduced standard search. Independent Q reconstruction
and proof reading likewise remain charged. Component queries, essentiality
and their proof costs are separate obligations. A small effective dimension
does not alone give a polynomial or quasi-polynomial total search when the
original Q phase has uncontrolled dimension. The general centre, distance
and hierarchy bounds required for unknot recognition remain unresolved.

\subsection{Source oracles, proof mutations and complete timings}
All 1,427 maintained tests pass in 289.263 seconds. The 28 focused tests
pass in 13.271 seconds. They check planar admission from raw nullity four,
generic reduction from five to four, complete parent records in the
original coordinate width, legacy replay without the support annotation,
shared-cap boundaries and interruption during source recompilation.
Eleven mutations remove or forge nested Q rays, change the source or allowed
support, change the nested phase or status, add recursive nesting, or alter
parent coverage and Euler data. They are all rejected with the producer,
Q iterator and planar producer disabled during independent replay.

The source audit repeats all 188 retained higher-nullity sectors against
release \code{c21e78bcd}. It completes in 53.634 seconds. Every complete
status agrees. For every source sector, complete standard rays from the
recompiled support equal the retained original-support source oracle in
actual full coordinates. Sixty-five sectors have forced-zero types: 22
fall from matching nullity four to three, 41 stay at four, and two fall from
five to four. Fresh Regina standard enumerations for all three source
triangulations involved reproduce the complete compatible ray lists in all
65 cases. These are actual finite source triangulations, not manufactured
post-elimination kernels.

Automatic discovery recompiles 61 of those sectors and returns 61 augmented
negative standard proofs. The other four return a verified essential Q
witness before recompilation is needed. All 22 raw-nullity-four sectors with
feasible nullity three enter the planar standard producer. Every augmented
parent proof retains the original allowed type list and coordinate width.
The unchanged original full-sector verifier also accepts representative
unannotated proofs at effective dimensions three and four. Independent
publication review replays all 61 augmented proofs with the producers
disabled, and also replays all 31 saved native diagram witnesses from the
preceding release.

The paired benchmark freezes both the preceding discovery module and its
independent checker. Each call constructs its source, completes standard
discovery and all configured work, replays its proof using that arm's
checker, and serializes the full result. Five randomized paired rounds
with identical-baseline A/A controls complete 100 measured calls and twenty
warmups. Both arms agree on every status; every proof is replayed and each
arm's certificate digest is stable across its repetitions.

\begin{center}
\small
\begin{tabular}{lrrrrr}
\toprule
Source & previous ms & current ms & old/new & old A/A & new A/A\\
\midrule
\code{cap\_3\_5\_2-4-to-3} & 164.057 & 41.699 & 3.863 & 1.045 & 1.005\\
\code{cap\_3\_5\_3-4-to-4} & 78.158 & 78.027 & 0.984 & 1.019 & 1.005\\
\code{cap\_3\_5\_3-5-to-4} & 3452.408 & 143.699 & 23.862 & 1.001 & 0.998\\
\code{cap\_1\_2\_3-4-to-4} & 40.810 & 40.753 & 1.004 & 1.002 & 1.002\\
\code{fibonacci\_lst\_09-4-to-4} & 60.635 & 61.058 & 0.999 & 0.988 & 0.988\\
\bottomrule
\end{tabular}
\end{center}


The planar-admission case has paired old/new ratio 3.863, with 684 versus
nine begun candidates, including its completed Q phase. Its incumbent A/A
ratio is 1.045 and the treatment A/A ratio is 1.005; the gain is substantially
larger than those fluctuations. The nullity-five-to-four case has ratio
23.862, with 12,655 versus 369 candidates and A/A controls about 1.001 and
0.998. Both the producer and independent replay avoid their respective
original larger standard arrangements.

The forced-zero case that stays at dimension four is near parity at 0.984;
its candidate count remains 124 in both arms. The full-support positive
control and nonpositive Q control have ratios 1.004 and 0.999. Removing
coordinates without changing effective dimension does not supply a uniform
speedup, and the larger nested proof and source recompilation are included
in the treatment cost. No source-query gain is attributed to improved
coverage of arbitrary knot diagrams.

The runtime release is \code{936e83365}. Tests, the complete source audit,
fresh Regina controls and paired timings are retained under
\code{data/feasible-span-*}. From \code{fast/}, use the module
\code{normal\_orbit\_research.feasible\_span} with its \code{audit} and
\code{benchmark} modes; the audit accepts \code{--fresh-regina}.
Publication checks verify all 609 runtime pins, 607 preceding-release
pins, every timing outcome, the complete source-sector verdicts and all
producer-disabled proof replays. The ray-geometry improvement and feasible
span theorem are local results. A general quasi-polynomial recognizer and
an implemented polynomial support producer remain separate unresolved
requirements. The retained LP formulation gives a route to the latter; this
Q-list integration does not implement that solver.

